% =============================================================================
% Benchmark — ClaudeCodeTrace dataset
% =============================================================================

\section{The \dataset{} Benchmark}
\label{sec:benchmark}

The evaluation in \S\ref{sec:evaluation} is only as credible as
the traffic it replays. Existing LLM-serving benchmarks are not a good
fit: ShareGPT-derived corpora capture single-shot conversations with no
tool-use structure; SWE-bench-Verified~\cite{swe-bench-verified,
swe-bench-iclr} provides tasks but not on-wire request bodies; and the
preference-style benchmarks (MT-Bench~\cite{mtbench},
Alpaca\-Eval~\cite{alpacaeval}, Chatbot
Arena~\cite{chatbot-arena}) are about model output ranking, not cache
behavior. None of them produce the long, system-prompt-heavy,
tool-call-laced request bodies that a real agent harness emits, and
none of them exercise the post-compaction prefix pattern that motivates
\sys{} (\S\ref{sec:intro}).

\dataset{} fills this gap. It is a corpus of \textbf{13 captures
spanning 182 turns and 411k total tokens (389k prompt~+~22k
completion)} of on-wire Claude Code request bodies, captured through
the \sys{} proxy against a real Claude Code client running real
workloads on a researcher's laptop. It is released CC-BY 4.0 on
Hugging Face~\cite{claudecode-trace-hf} so that future work on agent
serving --- whether on cacheblend, on subsequent KV-reuse mechanisms,
or on entirely different infrastructure --- can use the same baseline.

\subsection{Subset composition}
\label{sec:dataset-subsets}

The corpus is organized into three subsets, each targeting a different
slice of agent traffic:

\begin{table*}[t]
\small
\centering
\begin{tabular}{lrrrrr}
\toprule
Subset & Captures & Turns & Substantive turns & Prompt tokens & Completion tokens \\
\midrule
\texttt{swebench\_verified/} & 5 & 24 & 19 & 121{,}741 & 555 \\
\texttt{post\_compact/} & 7 & 105 & 98 & 186{,}464 & 16{,}803 \\
\texttt{skill\_invocation/} & 1 & 53 & 52 & 80{,}660 & 4{,}334 \\
\midrule
\textbf{Total} & \textbf{13} & \textbf{182} & \textbf{169} & \textbf{388{,}865} & \textbf{21{,}692} \\
\bottomrule
\end{tabular}
\caption{\dataset{} subset composition. ``Substantive turns'' counts turns with
prompt token count $\geq$256 --- the threshold below which
\texttt{lmcache}'s \texttt{blend\_min\_tokens} guard prevents cacheblend
from engaging. The 5 captures' completion-token totals on
\texttt{swebench\_verified} are small because most turns terminate in
tool-call JSON of 30--100 tokens; the totals on \texttt{post\_compact}
and \texttt{skill\_invocation} are larger because those captures include
prose continuations.}
\label{tab:dataset-summary}
\end{table*}

\paragraph{\texttt{swebench\_verified/}} Five Claude Code sessions, each
working on one task from SWE-bench Verified~\cite{swe-bench-verified}:
\texttt{astropy-13398}, \texttt{matplotlib-26113}, \texttt{pylint-7080},
\texttt{pytest-7490}, \texttt{sphinx-11510}. Each session is captured
through its first 5 turns: a planning turn, several file-exploration
tool calls (\texttt{Read}, \texttt{Bash}, \texttt{Grep}), and an initial
patch attempt. Prompts grow from $\approx$1k tokens on turn~1 to
$\geq$6k tokens by turn~5 as tool outputs accumulate in the running
conversation. This subset is the canonical baseline for agent
prefix-cache work: small-but-realistic, faithful to a published task
suite, and reproducible end-to-end (the original SWE-V task ID is
preserved in each capture's metadata).

\paragraph{\texttt{post\_compact/}} Seven captures (v1--v7) of CC's
\texttt{/compact} boundary, each running a single agentic conversation
through pre-compact buildup, an explicit \texttt{/compact}, and
post-compact continuation. \textbf{v1--v5 are diagnostic traces only ---
they were captured against mid-bug proxy states during development of
Patch~6 and the per-turn header normalizer (\S\ref{sec:patches}) and
are not part of this paper's benchmark claims.} v6 and v7 are the
canonical post-fix captures used in the evaluation. We expose all
seven in \dataset{} because the diagnostic ones are useful for
anyone debugging cache-LOOKUP behavior on Llama-class backends, but
the paper's headline numbers exclude v1--v5. Per-capture statistics are nearly uniform (each
runs the same 15-turn agent script): $\approx$26k prompt tokens, 14
substantive turns, max prompt $\approx$2.3k. The uniformity is
deliberate --- the subset is a controlled set of seven runs of the same
workload, varying only the proxy/cacheblend state, so any cross-capture
delta isolates a code change rather than a workload change.

\paragraph{\texttt{skill\_invocation/}} One capture (53 turns, 80k
prompt tokens, 52 substantive) that exercises CC's skill-loading
mechanism. We script CC to invoke five different skills (\texttt{bash-safety},
\texttt{commit-msg-style}, \texttt{markdown-table},
\texttt{python-import-order}, \texttt{weather-format}) across three
prompts each, plus a few setup and verification turns. The shared
agentic preamble across turns means cacheblend has substantial
within-capture prefix overlap to mine; the skill-specific suffixes
exercise the LOOKUP path with different selective-recompute patterns.
This subset is the smallest by capture count but the densest by turn
count, and it was the highest-value target for the statistical span
miner (\S\ref{sec:miner}) --- 49 of the top 100 mined spans come from
this single capture's prompt vocabulary.

\subsection{On-wire schema}
\label{sec:dataset-schema}

Each capture is one directory on disk with the following files:

\begin{itemize}
\item \texttt{traces.sqlite} --- per-request metadata: \texttt{request\_id},
  \texttt{session\_id}, \texttt{ts\_start}/\texttt{ts\_end},
  \texttt{prompt\_token\_count}, \texttt{response\_token\_count}, the
  full \texttt{request\_body\_json}, and instrumented cacheblend
  counters (\texttt{cache\_read\_tokens},
  \texttt{cache\_recompute\_tokens},
  \texttt{chunk\_aligned\_hit\_tokens},
  \texttt{tokens\_recomputed\_hkvd},
  \texttt{invariant\_violations}). These last five are populated when
  the capture ran against an instrumented lmcache build and surface the
  internal accounting that the §1.6 patches rely on; for plain proxy
  captures (e.g., \texttt{post\_compact/v1}) they are NULL.
\item \texttt{tokens/<request\_id>.parquet} --- the tokenized prompt as
  emitted to vLLM, one row per chunk (256 tokens by default), with the
  chunk hash. This is what cacheblend's LOOKUP key is computed against
  and what the statistical span miner walks.
\item \texttt{lookups/<request\_id>.parquet} --- per-chunk LOOKUP
  results when the capture's pod had cacheblend enabled (HIT / MISS /
  partial), with the chunk hash, the lookup latency, and the rescue
  contribution.
\item \texttt{vllm.log}, \texttt{oneshot\_boot.log},
  \texttt{proxy.log} --- per-pod logs, useful for cross-referencing
  the per-request rescue numbers against engine-side counters.
\end{itemize}

The schema is designed for additive evolution. New cacheblend
counters can be added as new \texttt{sqlite} columns without breaking
existing reader code; new per-chunk fields can be added to the parquet
schemas via schema-on-read. The \texttt{request\_body\_json} column is
the authoritative source of truth --- everything else can be
recomputed from it.

\subsection{Redaction surface}
\label{sec:dataset-redact}

CC's on-wire request body and our proxy's logs both incidentally
contain credentials and identifiers that must not leave the
researcher's laptop. We treat redaction as a release blocker, not an
afterthought.

\sys{} ships a redaction utility (\texttt{scripts/redact.py}) with
named patterns for five identifier classes:

\begin{enumerate}
\item \textbf{Tailscale} --- auth keys (\texttt{tskey-auth-*}) and host
  identifiers in URLs.
\item \textbf{Hugging Face} --- access tokens (\texttt{hf\_*}).
\item \textbf{RunPod} --- API keys (\texttt{rpa\_*}) and per-pod ingress
  URLs.
\item \textbf{Filesystem} --- absolute paths under \texttt{/Users/<user>/}
  rewritten to \texttt{/Users/<user>/}.
\item \textbf{CC client identifiers} --- the long-lived CC version
  hash and the user-account ID that appear in request headers.
\end{enumerate}

A second utility (\texttt{scripts/publish\_claudecode\_trace.py}) re-runs
the same patterns over every released file (\texttt{.json},
\texttt{.log}, \texttt{.txt}, \texttt{.md}, plus parquet/sqlite payloads
walked separately) and fails the publish on any post-redaction match
that is not itself a redaction marker. The audit reports
\textbf{0 violations} on the final 6.6\,MB tarball that was uploaded to
Hugging Face. The audit log itself is included in the dataset card as
provenance.

\subsection{License and access}
\label{sec:dataset-license}

\dataset{} is released as
\texttt{intelchen/claudecode-trace}~\cite{claudecode-trace-hf} under
CC-BY 4.0. The dataset card documents the schema, the redaction surface,
the SWE-V task IDs covered, and the
\texttt{LMCACHE\_BLEND\_RECOMPUTE\_RATIOS}=0.15 default that every
substantive measurement in \S\ref{sec:evaluation} uses. A reader
download script
(\texttt{scripts/download\_claudecode\_trace.py}) wraps Hugging Face's
\texttt{snapshot\_download} with subset filters so that downstream
benchmarks can pull only the captures they need.
